\section{Response-Time Analysis}

This section derives the response-time bounds for callbacks managed by RTExecutor.
Since the nested scheduling structure is eliminated---all callbacks execute as RTEMS tasks under a single fixed-priority preemptive scheduler---classic real-time scheduling theory becomes directly applicable.

\subsection{Scheduling Model}

We adopt the following assumptions:
\begin{enumerate}
    \item The RTEMS kernel employs fixed-priority preemptive scheduling with FIFO tie-breaking (FP-FIFO) on a single-core processor.
    \item Each ROS~2 callback $\gamma_i$ in the callback set $\Gamma = \{\gamma_1, \gamma_2, \dots, \gamma_n\}$ is mapped to a dedicated RTEMS task $\tau_i$ with priority $\pi_i$, period $T_i$, worst-case execution time $C_i$, and relative deadline $D_i$.
    \item Priorities are distinct and assigned according to the rate-monotonic policy: a shorter period implies a higher priority.
    \item RTEMS priority inheritance is enabled on all shared resources (e.g., the \texttt{RtsemQueue}), bounding priority inversion.
\end{enumerate}

Because RTExecutor delegates all scheduling decisions to the RTEMS kernel, the middleware-layer scheduler no longer exists as an independent scheduling entity.
Consequently, the response-time analysis degenerates to the well-studied classic fixed-priority preemptive scheduling model, and no Executor-level blocking terms need to be accounted for.

\subsection{Timer Callback Response Time}

\subsubsection{Release Model}

The TimeManager releases timer callbacks through the following chain:
\begin{equation}
    \text{HW tick interrupt} \;\to\; \text{Watchdog expiration} \;\to\; \text{Timer Server notification} \;\to\; \text{Worker execution}.
\end{equation}
At each clock tick, the Clock Tick ISR fires and inspects the RTEMS Watchdog red-black tree.
If a timer's expiration time has been reached, the Timer Server task is notified, which in turn signals the corresponding Worker task via an RTEMS event or semaphore.
The Worker task then executes the registered timer callback.

\subsubsection{Release Jitter}

The timer callback release suffers from jitter $J_i^{\mathrm{timer}}$ bounded by:
\begin{equation}
    J_i^{\mathrm{timer}} \leqslant 2 \; \text{ticks}.
\end{equation}
This bound arises from four sources:
\begin{itemize}
    \item \textbf{Tick quantization error} ($\leqslant 1$~tick): The hardware timer fires at the next tick boundary after the requested expiration time, which may be up to one tick early or late relative to the ideal release instant.
    \item \textbf{ISR execution delay}: The Clock Tick ISR must traverse the Watchdog red-black tree; the traversal time adds a bounded delay before the Timer Server is notified.
    \item \textbf{Timer Server context switch}: The Timer Server task must be dispatched by the scheduler before it can signal the Worker task.
    \item \textbf{Worker notification}: The \texttt{rtems\_event\_send} operation and the subsequent context switch to the Worker task contribute at most one additional tick of delay.
\end{itemize}
In our experimental platform (ARM \texttt{realview\_pbx\_a9\_qemu} with $\text{CONFIGURE\_MICROSECONDS\_PER\_TICK} = 1000$), the measured jitter lies in the range $[0,\,2]$~ticks (i.e., 0--2~ms), confirming the bound.

\subsubsection{Response-Time Bound}

Under the FP-FIFO scheduling model with release jitter, the worst-case response time $R_i$ of a timer callback $\gamma_i$ satisfies the following recurrent equation:
\begin{equation}
    R_i = C_i + \sum_{j \in hp(i)} \left\lceil \frac{R_i + J_j}{T_j} \right\rceil \times C_j + B_i,
    \label{eq:rta-timer}
\end{equation}
where $hp(i)$ denotes the set of callbacks with higher priority than $\gamma_i$, $J_j$ is the release jitter of callback $\gamma_j$, and $B_i$ is the worst-case blocking time due to priority inversion on shared resources protected by priority-inheritance semaphores.
The recurrence is solved iteratively starting from $R_i^{(0)} = C_i$ until $R_i^{(k+1)} = R_i^{(k)}$ or $R_i^{(k)} > D_i$ (unschedulable).

\subsection{Subscription Callback Response Time}

\subsubsection{Release Model}

The EventManager dispatches subscription callbacks through the following path:
\begin{equation}
    \texttt{rcl\_wait} \;\to\; \texttt{RtsemQueue.push()} \;\to\; \text{RTEMS event} \;\to\; \text{Worker execution}.
\end{equation}
The spin thread invokes \texttt{rcl\_wait} to detect subscription readiness.
Upon detection, the callback is wrapped as a \texttt{std::function<void()>} and pushed into the \texttt{RtsemQueue}.
An RTEMS event is then sent to the corresponding Worker task, which dequeues and executes the callback.

\subsubsection{Release Jitter}

The subscription callback release jitter is given by:
\begin{equation}
    J_i^{\mathrm{sub}} = t_{\mathrm{wait}} + t_{\mathrm{push}} + t_{\mathrm{ctx}},
\end{equation}
where $t_{\mathrm{wait}}$ is the delay from subscription readiness detection by \texttt{rcl\_wait}, $t_{\mathrm{push}}$ is the time to acquire the \texttt{RtsemQueue} semaphore and enqueue the callback, and $t_{\mathrm{ctx}}$ is the context-switch time from the spin thread to the Worker task.
Under RTEMS, both $t_{\mathrm{push}}$ and $t_{\mathrm{ctx}}$ are bounded and deterministic.

\subsubsection{End-to-End Latency}

For an intra-process publish--subscribe chain, the end-to-end latency from the publisher's callback invocation to the subscriber's callback completion is decomposed as:
\begin{equation}
    L_{\mathrm{e2e}} = L_{\mathrm{pub}} + L_{\mathrm{intra}} + L_{\mathrm{disp}} + L_{\mathrm{exec}},
    \label{eq:e2e}
\end{equation}
where:
\begin{itemize}
    \item $L_{\mathrm{pub}}$ is the publisher callback execution time (i.e., $C_{\mathrm{pub}}$).
    \item $L_{\mathrm{intra}}$ is the intra-process zero-copy delivery latency, which consists of the \texttt{rcl\_wait} detection delay and the data-copy cost. Under intra-process communication, no serialization or inter-process transport is involved.
    \item $L_{\mathrm{disp}}$ is the EventManager dispatch latency, bounded by one tick plus one context switch: $L_{\mathrm{disp}} \leqslant 1~\text{tick} + t_{\mathrm{ctx}}$.
    \item $L_{\mathrm{exec}}$ is the subscriber callback execution time (i.e., $C_{\mathrm{sub}}$).
\end{itemize}

\subsection{Comparison with Nested Scheduling Analysis}

Under the standard ROS~2 Executor, nested scheduling introduces an Executor-level blocking term $B^{\mathrm{exec}}$ into the response-time analysis~\cite{casini2019rtas,tang2020rtss,blass2021rtss,jiang2022rtss}.
Specifically, a high-priority callback may be blocked by a lower-priority callback that currently occupies the Executor thread.
The resulting response-time bound takes the form:
\begin{equation}
    R_i^{\mathrm{nested}} = C_i + \sum_{j \in hp(i)} \left\lceil \frac{R_i + J_j}{T_j} \right\rceil \times C_j + B_i + B_i^{\mathrm{exec}},
    \label{eq:rta-nested}
\end{equation}
where $B_i^{\mathrm{exec}}$ represents the worst-case Executor-level blocking.
This term is difficult to bound tightly: when callbacks within the same \texttt{MutuallyExclusive} callback group share an Executor thread, $B_i^{\mathrm{exec}}$ can be as large as the longest execution time of any lower-priority callback in the group~\cite{casini2019rtas}.
In the worst case, when the Executor lacks priority distinction among callbacks, $B_i^{\mathrm{exec}}$ is effectively unbounded.

With RTExecutor, nested scheduling is eliminated entirely.
The Executor-level blocking term vanishes:
\begin{equation}
    B_i^{\mathrm{exec}} = 0.
\end{equation}
Equation~\eqref{eq:rta-timer} therefore degenerates to the classic FP scheduling response-time analysis, which is both simpler and tighter than the nested analysis of Eq.~\eqref{eq:rta-nested}.
This constitutes a fundamental analytical advantage of the proposed approach: by collapsing the two-level scheduling hierarchy into a single OS-level scheduler, the worst-case response-time bounds become both easier to compute and more favorable for schedulability.

\subsection{Schedulability Condition}

A callback set $\Gamma$ is deemed schedulable under RTExecutor if and only if every callback meets its deadline:
\begin{equation}
    \forall \gamma_i \in \Gamma : \quad R_i \leqslant D_i,
    \label{eq:sched}
\end{equation}
where $R_i$ is computed via Eq.~\eqref{eq:rta-timer}.

When priorities are assigned according to the rate-monotonic policy ($D_i = T_i$), Liu and Layland's utilization test provides a sufficient schedulability condition:
\begin{equation}
    U = \sum_{i=1}^{n} \frac{C_i}{T_i} \leqslant n\left(2^{1/n} - 1\right).
    \label{eq:util}
\end{equation}
Note that Eq.~\eqref{eq:util} does not account for release jitter.
When jitter is non-negligible (e.g., for timer callbacks with $J_i^{\mathrm{timer}} \leqslant 2$~ticks), the exact response-time test of Eq.~\eqref{eq:rta-timer} should be applied instead.
Nevertheless, the key observation is that this analysis is \emph{identical} to the standard uniprocessor FP scheduling analysis---no Executor-specific terms need to be introduced, which is the direct consequence of eliminating nested scheduling.
